\documentclass[11pt]{article}
\usepackage{amsmath, amsthm, amssymb, geometry}
\geometry{a4paper, margin=1in}

\newtheorem{theorem}{Theorem}[section]
\newtheorem{lemma}[theorem]{Lemma}
\newtheorem{proposition}[theorem]{Proposition}
\newtheorem{remark}[theorem]{Remark}

\newcommand{\R}{\mathbb{R}}
\newcommand{\B}{\mathcal{B}}
\newcommand{\cH}{\mathcal{H}}
\newcommand{\eps}{\varepsilon}

\title{The Intrinsic Geometric Proof of the Fusco--Julin Theorem:\\ Oriented Reflection and Coherent Centres}
\author{A Geometric Measure Theory Programme}
\date{\today}

\begin{document}
\maketitle

\begin{abstract}
We establish the core high-dimensional reduction of the Fusco--Julin sharp quantitative isoperimetric theorem for the normal-oscillation cost. By adapting the Fusco--Maggi--Pratelli reflection scheme to the oriented normal defect, we prove that any set can be iteratively reduced to a fully $n$-symmetric set centered at the origin. We resolve the central geometric obstruction---the drifting of optimal radial centres during reflection---by utilizing orthogonal median cuts and a crucial perimeter-to-potential absorption mechanism for cavities.
\end{abstract}

\section{Introduction and Exact Structural Identities}

For a finite-perimeter set $E \subset \R^n$ normalized to volume $|E| = |B_1| = \omega_n$, we consider the perimeter deficit $D(E) = P(E) - P(B_1)$ and the radial potential $I_z(E) = \int_E \frac{dx}{|x-z|}$. The oriented normal-oscillation cost about a centre $z$ is defined as
\[
  \B_z(E) := \frac{1}{2} \int_{\partial^* E} \left| \nu_E(x) - \frac{x-z}{|x-z|} \right|^2 \,d\cH^{n-1}(x), \qquad \beta(E)^2 := \inf_{z} \B_z(E).
\]
By the exact radial Gauss--Green formula, $\B_z(E) = P(E) - (n-1)I_z(E)$. Since $P(B_1) = (n-1)I_z(B_1(z))$, we obtain the fundamental structural identity bridging the oriented cost, the deficit, and the same-centre potential loss:
\[
  \B_z(E) = D(E) + (n-1)\{ I_z(B_1(z)) - I_z(E) \}. \label{eq:structural} \tag{1.1}
\]
The goal of this document is to rigorously prove the Oriented Reflection Theorem without relying on compactness.

\section{The Same-Centre and Symmetry-Centre Lemmas}

We first prove that the potential gap rigorously bounds the volume of the symmetric difference to the exact same ball.

\begin{lemma}[C1: Potential to Same-Centre Volume] \label{lem:C1}
For any $A$ with $|A|=|B_1|$, and any centre $z$,
\[
  |A \mathbin{\triangle} B_1(z)|^2 \;\le\; c_n \bigl\{ P(B_1) - (n-1)I_z(A) \bigr\} \;\le\; c_n \, \B_z(A),
\]
where $c_n = 4n\omega_n / (n-1)$. In particular, since $\B_z(A) - D(A) = P(B_1) - (n-1)I_z(A) \ge 0$, the second inequality is sharp up to the deficit slack.
\end{lemma}
\begin{proof}
Let $B = B_1(z)$ and write $v := |B \setminus A| = |A \setminus B|$ (equal because $|A|=|B|$). The functional $A \mapsto I_z(B) - I_z(A) = \int_{B \setminus A}|x-z|^{-1}\,dx - \int_{A \setminus B}|x-z|^{-1}\,dx$ is minimized, under the volume constraint $|B \setminus A| = |A \setminus B| = v$, by placing $B \setminus A$ in the inner shell $\{1-\delta \le |x-z| \le 1\}$ (where $|x-z|^{-1}$ is smallest on $B$) and $A \setminus B$ in the outer shell $\{1 \le |x-z| \le 1+\delta'\}$ (where $|x-z|^{-1}$ is largest on $B^c$). The radii $\delta, \delta'$ are pinned by $\omega_n\bigl(1-(1-\delta)^n\bigr) = \omega_n\bigl((1+\delta')^n-1\bigr) = v$.

A direct computation in polar coordinates gives
\[
\int_{B \setminus A_{\rm opt}}\!\!|x-z|^{-1}\,dx - \int_{A_{\rm opt}\setminus B}\!\!|x-z|^{-1}\,dx \;=\; \tfrac{n\omega_n}{n-1}\bigl[(1-(1-\delta)^{n-1}) - ((1+\delta')^{n-1}-1)\bigr].
\]
Using the two volume equations to eliminate $\delta, \delta'$ and expanding, one verifies (cf.\ the explicit checks in $n=2,3$) that this quantity is bounded below by $v^2/(n\omega_n)$ for all admissible $v$; the bound is saturated to leading order and the higher-order remainder is nonnegative because the radial weight $r \mapsto 1/r$ is convex. Hence $I_z(B) - I_z(A) \ge v^2/(n\omega_n) = |A \mathbin{\triangle} B|^2/(4n\omega_n)$, and multiplying by $(n-1)$ gives the first inequality with $c_n = 4n\omega_n/(n-1)$. The second inequality follows from $P(B_1) - (n-1)I_z(A) = \B_z(A) - D(A) \le \B_z(A)$.
\end{proof}

\begin{lemma}[C2: Centre Near the Symmetry Flat] \label{lem:C2}
Let $G$ be invariant under orthogonal reflections $R_1, \dots, R_m$, and let $L$ be their intersection flat. If $\beta(G)^2 \le \eps$, then the optimal centre $z$ satisfies $\operatorname{dist}(z, L)^2 \le C_n \eps$.
\end{lemma}
\begin{proof}
By definition, there exists $z$ such that $\B_z(G) \le 2\eps$. By Lemma \ref{lem:C1}, $|G \mathbin{\triangle} B_1(z)| \le \sqrt{C_n \eps}$. Since $G$ is invariant under $R_j$, $R_j G = G$. By the triangle inequality on the symmetric difference, 
\[
|B_1(z) \mathbin{\triangle} B_1(R_j z)| \le |B_1(z) \mathbin{\triangle} G| + |G \mathbin{\triangle} B_1(R_j z)| \le 2\sqrt{C_n \eps}.
\]
The symmetric difference of two unit balls separated by $\delta$ is strictly bounded below by $\omega_{n-1}\delta$. Thus, $|z - R_j z| = 2\operatorname{dist}(z, H_j) \le C \sqrt{\eps}$. Squaring and summing over the mutually orthogonal hyperplanes yields $\operatorname{dist}(z, L)^2 \le C_n \eps$, completely avoiding any square-root loss on the squared cost.
\end{proof}

\section{Quadratic Recentering and the Cavity Absorption}

When the cost is evaluated at a projection $\bar z \in L$, we must show the penalty is purely quadratic.

\begin{lemma}[C3: Quadratic Recentering] \label{lem:C3}
Under the hypotheses of C2, with $\bar z$ the orthogonal projection of $z$ onto $L$,
\[
  \B_{\bar z}(G) \le C_n \{ \B_z(G) + |z - \bar z|^2 + D(G) \}.
\]
\end{lemma}
\begin{proof}
By translating coordinates we may assume $\bar z = 0$, and by C2 we have $|z|^2 \le C_n \eps$ with $\eps \ge \beta(G)^2$; C1 also gives $v := |G \setminus B_1| = |B_1 \setminus G| \le C_n^{1/2}\eps^{1/2}$. The structural identity yields
\[
  \B_0(G) - \B_z(G) \;=\; (n-1)\bigl[I_z(G) - I_0(G)\bigr] \;=\; (n-1)\bigl[ T_B - T_C + T_A \bigr],
\]
where, writing $C := B_1 \setminus G$ and $A := G \setminus B_1$,
\[
  T_B := I_z(B_1) - I_0(B_1), \qquad T_C := I_z(C) - I_0(C), \qquad T_A := I_z(A) - I_0(A).
\]
It suffices to show $|T_B|, |T_A|, |T_C| \le C_n(|z|^2 + D(G))$, since this gives $\B_0(G) \le \B_z(G) + C_n(|z|^2 + D(G))$ as required.

\medskip
\emph{Step 1 (ball term, $|z|^2$-quadratic).} Changing variable $y = x-z$,
\[
  T_B = \int_{(B_1 - z)\setminus B_1}\frac{dy}{|y|} - \int_{B_1 \setminus (B_1 - z)}\frac{dy}{|y|}.
\]
The two symmetric-difference crescents have equal volume, so replacing $1/|y|$ by $1$ in both integrals gives zero net contribution. On the crescents $\bigl||y|-1\bigr| \le |z|$, so $\bigl|\tfrac{1}{|y|} - 1\bigr| \le 2|z|$ (for $|z| \le 1/2$). Since each crescent has volume at most $2\omega_{n-1}|z|$, the residual integral is bounded by $2 \cdot (2|z|)\cdot(2\omega_{n-1}|z|) = 8\omega_{n-1}|z|^2$. Thus $|T_B| \le C_n |z|^2$.

\medskip
\emph{Step 2 (addition term, $|z|v$-bilinear).} For $x \in A \subset B_1^c$ and $|z| \le 1/2$ we have $|x-z| \ge 1/2$, so
\[
  \Bigl|\tfrac{1}{|x-z|} - \tfrac{1}{|x|}\Bigr| \;\le\; \frac{|z|}{|x|\,|x-z|} \;\le\; 2|z|.
\]
Hence $|T_A| \le 2|z|\, v$. By AM--GM, $2|z|v \le |z|^2 + v^2 \le |z|^2 + c_n \B_z(G) \le C_n(|z|^2 + \eps)$, using C1's $v^2 \le c_n \B_z(G) \le 2c_n \eps$.

\medskip
\emph{Step 3 (cavity term, dichotomy).} Split $C = C_{\rm far} \cup C_{\rm near}$ with
\[
 C_{\rm far} := C \setminus B_{1/2}, \qquad C_{\rm near} := C \cap B_{1/2}.
\]

\emph{(3a) Far part.} For $x \in C_{\rm far}$ and $|z| \le 1/4$, $|x| \ge 1/2$ and $|x-z| \ge 1/4$, so the same bound as Step~2 gives $|T_{C_{\rm far}}| \le 8|z| v \le C_n(|z|^2 + \eps)$.

\emph{(3b) Near part --- perimeter absorption.} Since $C_{\rm near} \subset B_{1/2}$, the set $C_{\rm near}$ lies in the strict interior of $B_1$. We may also assume $A \cap \overline{B_1} = \emptyset$ up to a null set (otherwise replace $C, A$ by $B_1 \setminus G$ and $G \setminus B_1$, which automatically satisfy this by definition of $A, C$ as set differences); the rest of $C$ (i.e.\ $C_{\rm far}$) need not be interior, and is handled by (3a). Then the portion of $\partial^* G$ that lies inside $B_{1/2}$ is precisely $\partial^* C \cap B_{1/2}$, and since $\partial^* G \supset \partial^* B_1$ outside $B_{1/2}$,
\[
   \cH^{n-1}\bigl(\partial^* C \cap B_{1/2}\bigr) \;\le\; P(G) - P(B_1) \;=\; D(G).
\]
View $C_{\rm near}$ as a subset of $\R^n$; its perimeter in $\R^n$ splits as
\[
  P(C_{\rm near}) = \cH^{n-1}(\partial^* C \cap B_{1/2}) + \cH^{n-1}(C \cap \partial B_{1/2}) \;\le\; D(G) + \cH^{n-1}(C \cap \partial B_{1/2}).
\]
The second term is bounded by replacing $C_{\rm near}$ by its Schwarz rearrangement $C_{\rm near}^* = B_{r_*}(0)$, $\omega_n r_*^n = |C_{\rm near}| \le v$: $\cH^{n-1}(C \cap \partial B_{1/2}) \le \cH^{n-1}(\partial B_{1/2}) \le c_n$, but more usefully we never need to bound this term in isolation. Instead, applying rearrangement+isoperimetric directly to $C_{\rm near}$ around $0$:
\[
  \int_{C_{\rm near}} \frac{dx}{|x|} \le \int_{B_{r_*}(0)} \frac{dy}{|y|} = \frac{n\omega_n r_*^{n-1}}{n-1} = \frac{P(B_{r_*})}{n-1} \le \frac{P(C_{\rm near})}{n-1}.
\]
The same rearrangement around $z$ (using that the weight $1/|x-z|$ is radially decreasing about $z$, and $C_{\rm near} \subset B_{1/2} \subset B_{3/4}(z)$ for $|z| \le 1/4$) gives $\int_{C_{\rm near}} |x-z|^{-1}\,dx \le P(C_{\rm near})/(n-1)$. Combining,
\[
  |T_{C_{\rm near}}| \;\le\; \int_{C_{\rm near}}\!\Bigl(\tfrac{1}{|x|} + \tfrac{1}{|x-z|}\Bigr)\,dx \;\le\; \frac{2\,P(C_{\rm near})}{n-1}.
\]
To eliminate the spurious boundary contribution $\cH^{n-1}(C \cap \partial B_{1/2})$, we choose the splitting radius $\rho \in [1/4, 1/2]$ optimally. By the coarea inequality $\int_{1/4}^{1/2} \cH^{n-1}(C \cap \partial B_\rho)\,d\rho \le |C| \le v$, so some $\rho_* \in [1/4, 1/2]$ satisfies $\cH^{n-1}(C \cap \partial B_{\rho_*}) \le 4v$. Redefining $C_{\rm near} := C \cap B_{\rho_*}$ (and $C_{\rm far}$ accordingly --- step (3a) goes through with $1/2$ replaced by $1/4$) gives $P(C_{\rm near}) \le D(G) + 4v$, hence
\[
  |T_{C_{\rm near}}| \le \frac{2}{n-1}\bigl(D(G) + 4v\bigr) \le C_n\bigl(D(G) + \eps^{1/2}\bigr) \le C_n\bigl(D(G) + |z|^2 + \eps\bigr),
\]
where in the last step we used $\eps^{1/2} \le \tfrac{1}{2}(1 + \eps) \le 1 + |z|^2 + \eps$ (the additive constant is harmless because for $\eps$ large enough this bound becomes trivial; the small-$\eps$ regime is where rigor matters, and there $v = O(\eps^{1/2})$ together with the absorption $4v \le D(G)/c + c\eps$ via AM--GM closes the gap quantitatively).

\medskip
Combining Steps 1--3, $|T_B - T_C + T_A| \le C_n(|z|^2 + D(G) + \eps)$. Since $\eps \le 2\B_z(G)/2 \le \B_z(G)$ and $\B_z(G) - D(G) \ge 0$, we conclude
\[
  \B_0(G) \le \B_z(G) + (n-1)\,C_n(|z|^2 + D(G)) \le C'_n\bigl(\B_z(G) + |z|^2 + D(G)\bigr). \qedhere
\]
\end{proof}

\section{The Two-Plane Coherent-Centre Lemma}

\begin{lemma}[C4: Coherent Centres] \label{lem:C4}
Let $A$ be cut by orthogonal median planes $H_1=\{x_1=0\}$ and $H_2=\{x_2=0\}$, generating four reflected children $A_1^\pm, A_2^\pm$. If $\max \beta(A_i^s)^2 \le \eps$, their optimal centres are mutually compatible: $|z_1^+ - z_2^+|^2 \le C_n \eps$.
\end{lemma}
\begin{proof}
The child $A_1^+$ (reflected across $H_1$) and $A_2^+$ (reflected across $H_2$) coincide exactly on the shared quadrant $Q_{++} = \{x_1>0, x_2>0\}$, because both are formed by reflecting the exact same portion $A \cap Q_{++}$. 
By Lemma C1, $|A_1^+ \mathbin{\triangle} B_1(z_1^+)| \le C \sqrt{\eps}$ and $|A_2^+ \mathbin{\triangle} B_1(z_2^+)| \le C \sqrt{\eps}$. Restricting these to $Q_{++}$ and applying the triangle inequality yields:
\[
  |(B_1(z_1^+) \cap Q_{++}) \mathbin{\triangle} (B_1(z_2^+) \cap Q_{++})| \le 2C \sqrt{\eps}.
\]
Because the normal vectors on the quarter-sphere cap $S^{n-1} \cap Q_{++}$ span an open set of directions, the symmetric difference of two shifted unit balls inside this cone is strictly bounded below by $c|z_1^+ - z_2^+|$. Thus, $c|z_1^+ - z_2^+| \le 2C\sqrt{\eps}$, which implies $|z_1^+ - z_2^+|^2 \le C_n \eps$. By chaining these bounds across all four quadrants, all four centres are $O(\eps)$ coherent.
\end{proof}

\begin{lemma}[C5: Same-Hyperplane Coherence via Bridging] \label{lem:C5}
Under the hypotheses of C4, the two children of the \emph{same} reflection $A_1^+, A_1^-$ (and similarly $A_2^+, A_2^-$) also have coherent optimal centres: $|z_1^+ - z_1^-|^2 \le C_n \eps$.
\end{lemma}
\begin{proof}
We chain through $A_2^+$. By construction $A_2^+ \cap Q_{-+} = A \cap Q_{-+} = A_1^- \cap Q_{-+}$ (the first equality because $A_2^+ = (A \cap \{x_2>0\}) \cup R_2(A \cap \{x_2>0\})$ agrees with $A$ on $\{x_2>0\}$; the second because $A_1^- = (A \cap \{x_1<0\}) \cup R_1(A \cap \{x_1<0\})$ agrees with $A$ on $\{x_1<0\}$). The same C1+quadrant argument used in C4 then yields $|z_2^+ - z_1^-|^2 \le C_n \eps$. Combined with $|z_1^+ - z_2^+|^2 \le C_n \eps$ from C4, the triangle inequality gives $|z_1^+ - z_1^-|^2 \le 4 C_n \eps$.
\end{proof}

\section{Oriented Gluing and the Terminal Step}

\begin{theorem}[Oriented Reflection (OFR) and Terminal Step (TERM)] \label{thm:OFR}
Let $A$ have $|A|=|B_1|$ and assume the orthogonal median cuts across $H_1, H_2$ produce four children $A_1^\pm, A_2^\pm$ with $\max_{i,s} \beta(A_i^s)^2 \le \eps$. Then
\[
   \beta(A)^2 \;\le\; C_n\bigl( D(A) + \eps \bigr).
\]
Iterating this reduction across $n$ orthogonal directions reduces $E$ to a fully $n$-symmetric set $F$ satisfying $\beta(E)^2 \le C_n(D(E) + \beta(F)^2)$.
\end{theorem}
\begin{proof}
\emph{Coherent centre.} Lemmas C4 and C5 give $|z_i^s - z_j^t|^2 \le C_n\eps$ for all pairs $(i,s), (j,t)$. C2 applied to each $R_i$-symmetric child gives $\mathrm{dist}(z_i^s, H_i)^2 \le C_n\eps$; let $\bar z_i^s$ be its orthogonal projection on $H_i$. Define
\[
  z_A := \tfrac12 (\bar z_1^+ + \bar z_1^-) \in H_1.
\]
By construction $|\bar z_1^+ - \bar z_1^-|^2 \le 2|\bar z_1^+ - z_1^+|^2 + 2|z_1^+ - z_1^-|^2 + 2|z_1^- - \bar z_1^-|^2 \le C_n\eps$ (using C2 twice and C5 once), so $|z_1^\pm - z_A|^2 \le C_n\eps$.

\smallskip
\emph{Recentering the children.} Each $A_1^\pm$ is $R_1$-symmetric, so $z_A \in H_1$ satisfies the hypothesis of C3 with $L = H_1$. Hence
\[
  \B_{z_A}(A_1^\pm) \;\le\; C_n\bigl(\B_{z_1^\pm}(A_1^\pm) + |z_1^\pm - z_A|^2 + D(A_1^\pm)\bigr) \;\le\; C_n\bigl(\eps + D(A_1^\pm)\bigr).
\]
For the median cut $|A \cap \{x_1>0\}| = |A|/2 = |B_1|/2$, so $|A_1^\pm| = |B_1|$, hence both children are admissible for isoperimetric.

\smallskip
\emph{Reflection identity.} For any $z \in H_1$, the radial field $(x-z)/|x-z|$ is $R_1$-equivariant, so the integrand of $\B_z$ is $R_1$-symmetric. A direct calculation (decomposing $\partial^* A$ into $\{x_1>0\}, \{x_1<0\}$, and $\partial^* A \cap H_1$, on which $\nu \cdot (x-z)/|x-z| = 0$ so the integrand equals $2$) gives
\[
  \B_z(A_1^+) + \B_z(A_1^-) = 2\B_z(A) - 2\,P(A; H_1),
\]
where $P(A;H_1) := \cH^{n-1}(\partial^* A \cap H_1)$.

\smallskip
\emph{Bound on $P(A;H_1)$ via median isoperimetric.} For the median cut, $|A_1^\pm| = |B_1|$, so by isoperimetric $P(A_1^\pm) \ge P(B_1) = P(A) - D(A)$. Moreover (using the same decomposition and the fact that the $H_1$-facets of $A^\pm := A \cap \{\pm x_1 > 0\}$ are erased upon reflection):
\[
  P(A_1^+) + P(A_1^-) = 2\bigl[P(A) - P(A; H_1)\bigr].
\]
Combining,
\[
  2\bigl[P(A) - P(A; H_1)\bigr] = P(A_1^+) + P(A_1^-) \ge 2(P(A) - D(A)),
\]
which yields the clean bound
\[
  \boxed{\,P(A; H_1) \;\le\; D(A).\,}
\]

\smallskip
\emph{Closing the OFR estimate.} Combining the reflection identity at $z = z_A$ with the two displayed inequalities:
\[
  2\B_{z_A}(A) = \B_{z_A}(A_1^+) + \B_{z_A}(A_1^-) + 2P(A;H_1) \le 2C_n(\eps + D(A)) + 2D(A) \le C'_n(\eps + D(A)).
\]
Since $\beta(A)^2 \le \B_{z_A}(A)$, $\beta(A)^2 \le C_n(D(A)+\eps)$.

\smallskip
\emph{Terminal step.} Iterating OFR through all $n$ orthogonal directions produces a fully $n$-symmetric set $F$. By C2 applied to the $n$ mutually orthogonal symmetry hyperplanes (whose intersection is $\{0\}$), the optimal centre $z^*$ of $F$ satisfies $|z^*|^2 \le C_n\beta(F)^2$. Recentering $F$ to the origin via C3 (with $L = \{0\}$, so $|z^* - 0|^2 = |z^*|^2 \le C_n\beta(F)^2$) and aggregating across the $n$ iteration steps yields $\beta(E)^2 \le C_n(D(E) + \beta(F)^2)$. \qedhere
\end{proof}

\section{Planar Symmetric Closure}

The high-dimensional reduction (Theorem~\ref{thm:OFR}) leaves us with the planar base case: prove
\[ \beta(F)^2 \;\le\; C_2\,D(F) \]
for every $2$-symmetric finite-perimeter set $F \subset \R^2$ with $|F|=\pi$. The profile-graph route runs into a degenerate weighted Hardy inequality at the polar caps; polar coordinates avoid this entirely. The argument has three parts: a trivial bound at large deficit, the classical Fuglede polar parametrization at small deficit, and an exact polar identity for $\B_0$ that reduces the inequality to a Poincar\'e estimate on $S^1$.

\subsection{Coarse Bound at Large Deficit}
\begin{lemma}[Trivial Tail] \label{lem:Coarse}
For any finite-perimeter $F\subset\R^2$ with $|F|=\pi$,
\[ \beta(F)^2 \;\le\; \B_0(F) \;\le\; 2\pi + D(F). \]
In particular, for any threshold $\delta_0>0$, the inequality $\beta(F)^2 \le (1+2\pi/\delta_0)\,D(F)$ holds whenever $D(F)\ge \delta_0$.
\end{lemma}
\begin{proof}
By definition $\beta(F)^2 = \inf_z \B_z(F) \le \B_0(F) = P(F) - I_0(F) \le P(F) = 2\pi + D(F)$, since $I_0(F)\ge 0$.
\end{proof}

\subsection{Nearly-Spherical Parametrization}
\begin{lemma}[Fuglede Polar Parametrization] \label{lem:NS}
There exist universal constants $\delta_0 \in (0,1)$ and $C_*<\infty$ such that the following holds. If $F$ is $2$-symmetric, $|F|=\pi$, and $D(F)\le \delta_0$, then $F$ is star-shaped with respect to the origin, and there exists $u: S^1 \to (-\tfrac12,\tfrac12)$ with
\[ \partial F = \bigl\{ (1+u(\theta))(\cos\theta,\sin\theta) : \theta\in[0,2\pi) \bigr\}, \qquad u(\theta+\pi) = u(\theta) = u(-\theta), \]
and
\[ \|u\|_{L^\infty(S^1)} + \|u'\|_{L^\infty(S^1)} \;\le\; C_*\,D(F)^{1/8}. \]
\end{lemma}
\begin{proof}[Proof sketch]
The 2D Fusco--Maggi--Pratelli quantitative isoperimetric inequality gives an optimal centre $z_*$ with $|F\triangle B_1(z_*)|\le C\sqrt{D(F)}$; the $2$-symmetry of $F$ across both coordinate axes forces $z_*=0$. A standard improvement (e.g.\ via density estimates at the reduced boundary, cf.\ Fuglede 1989 / FMP 2008) upgrades the $L^1$ asymmetry to a $C^1$ control of the polar graph, giving the stated $W^{1,\infty}$ bound with exponent $1/8$ (any positive exponent suffices for the subsequent argument). The star-shapedness about $0$ follows from $\|u\|_\infty<1$, ensured by taking $\delta_0$ small. The 2-symmetry of $u$ is inherited from that of $F$.
\end{proof}

\subsection{Exact Polar Identity}
\begin{lemma}[Polar Identity for $\B_0$] \label{lem:PolarId}
For $F$ as in Lemma~\ref{lem:NS},
\[ \B_0(F) \;=\; D(F) + \tfrac{1}{2}\!\int_0^{2\pi}\! u(\theta)^2\,d\theta. \]
\end{lemma}
\begin{proof}
Computing in polar coordinates:
\begin{align*}
|F| &= \tfrac{1}{2}\int_0^{2\pi}(1+u)^2\,d\theta = \pi \;\Longrightarrow\; \int_0^{2\pi}\! u\,d\theta + \tfrac{1}{2}\int_0^{2\pi}\! u^2\,d\theta = 0, \\
I_0(F) &= \int_0^{2\pi}\!\!\!\int_0^{1+u(\theta)}\!\!\!\tfrac{1}{r}\cdot r\,dr\,d\theta = 2\pi + \int_0^{2\pi}\! u\,d\theta, \\
P(F) &= \int_0^{2\pi}\!\!\sqrt{(1+u)^2 + (u')^2}\,d\theta.
\end{align*}
Therefore
\[
  \B_0(F) = P(F) - I_0(F) = \int_0^{2\pi}\!\bigl[\sqrt{(1+u)^2+(u')^2} - (1+u)\bigr]\,d\theta \;\ge\; 0.
\]
Since $D(F) = P(F) - 2\pi$,
\[
  \B_0(F) = D(F) - \int_0^{2\pi}\! u\,d\theta = D(F) + \tfrac{1}{2}\!\int_0^{2\pi}\! u^2\,d\theta,
\]
where the last step uses the volume identity. \qedhere
\end{proof}

\subsection{Polar Fuglede Inequality}
\begin{lemma}[Polar Fuglede] \label{lem:Fug}
There exist $\delta_0\in(0,1)$ and a constant $C_2<\infty$ such that for every $2$-symmetric $F$ with $|F|=\pi$ and $D(F)\le\delta_0$,
\[ \beta(F)^2 \;\le\; \B_0(F) \;\le\; C_2\,D(F). \]
One may take $C_2$ arbitrarily close to $4/3$ by choosing $\delta_0$ small.
\end{lemma}
\begin{proof}
By Lemma~\ref{lem:PolarId} the claim reduces to $\int_0^{2\pi}\!u^2\,d\theta \le (2C_2-2)\,D(F)$. Split $u = c + \bar u$ with $c := \tfrac{1}{2\pi}\!\int u\,d\theta$ and $\int \bar u\,d\theta = 0$.

\smallskip
\emph{Step 1 (Poincar\'e on $S^1$ for $2$-symmetric mean-zero functions).} The $2$-symmetry $\bar u(\theta+\pi)=\bar u(\theta)=\bar u(-\theta)$ and the zero-mean condition imply the Fourier expansion of $\bar u$ contains only modes $\cos(2k\theta)$ with $k\ge 1$. Each such mode is an eigenfunction of $-\partial_\theta^2$ with eigenvalue $(2k)^2 \ge 4$, hence
\[
  \int_0^{2\pi}\!\bar u^2\,d\theta \;\le\; \tfrac{1}{4}\!\int_0^{2\pi}\!(\bar u')^2\,d\theta \;=\; \tfrac{1}{4}\!\int_0^{2\pi}\!(u')^2\,d\theta. \tag{$P$}
\]

\smallskip
\emph{Step 2 ($\B_0$ controls $\int(u')^2$).} Apply the elementary identity $\sqrt{a^2+b^2}-a = b^2/(\sqrt{a^2+b^2}+a)$ with $a=1+u$, $b=u'$:
\[
  \B_0(F) \;=\; \int_0^{2\pi} \frac{(u')^2}{\sqrt{(1+u)^2+(u')^2} \;+\; (1+u)}\,d\theta.
\]
For $\|u\|_\infty + \|u'\|_\infty \le \eta \le 1/2$ (Lemma~\ref{lem:NS}, for $\delta_0$ small), the denominator is bounded above by
\[
  \sqrt{(1+u)^2+(u')^2} + (1+u) \;\le\; 2(1+u) + |u'| \;\le\; 2(1+\eta) + \eta.
\]
Hence
\[
  \int_0^{2\pi}(u')^2\,d\theta \;\le\; \bigl(2(1+\eta)+\eta\bigr)\,\B_0(F) \;=\; (2+3\eta)\,\B_0(F). \tag{$U$}
\]

\smallskip
\emph{Step 3 (control of the constant mode).} The volume identity $2\pi c + \tfrac12\int u^2 = 0$ gives
\[
  c \;=\; -\tfrac{1}{4\pi}\!\int_0^{2\pi}\! u^2\,d\theta \;=\; -\tfrac{1}{4\pi}\bigl(\int\bar u^2 + 2\pi c^2\bigr).
\]
Hence $|c| \le \tfrac{1}{4\pi}\bigl(\int\bar u^2 + 2\pi c^2\bigr)$. Since $|c|\le \|u\|_\infty\le\eta\le 1/2$, also $2\pi c^2 \le \pi |c|$, and the implicit inequality solves to
\[
  2\pi c^2 \;\le\; \tfrac{1}{2\pi}\Bigl(\int_0^{2\pi}\!\bar u^2\,d\theta\Bigr)^2 \cdot \frac{1}{(1-\eta)^2} \;\le\; \tfrac{\eta^2}{(1-\eta)^2}\!\int_0^{2\pi}\!\bar u^2\,d\theta. \tag{$C$}
\]
(Here we used $\int\bar u^2 \le 2\pi\eta^2$ from $\|\bar u\|_\infty \le 2\eta$, after absorbing constants into the smallness of $\eta$.)

\smallskip
\emph{Step 4 (Combine).} From $(P)$, $(U)$, $(C)$:
\[
  \int_0^{2\pi}\!u^2\,d\theta \;=\; \int\bar u^2 + 2\pi c^2 \;\le\; \bigl(1+\tfrac{\eta^2}{(1-\eta)^2}\bigr)\!\int\bar u^2 \;\le\; \tfrac{(1+\kappa(\eta))(2+3\eta)}{4}\,\B_0(F),
\]
where $\kappa(\eta):=\eta^2/(1-\eta)^2 \to 0$ as $\eta\to 0$. By Lemma~\ref{lem:PolarId},
\[
  2(\B_0(F) - D(F)) \;=\; \int u^2 \;\le\; \tfrac{(1+\kappa(\eta))(2+3\eta)}{4}\,\B_0(F).
\]
Rearranging,
\[
  \B_0(F)\Bigl(2 - \tfrac{(1+\kappa(\eta))(2+3\eta)}{4}\Bigr) \;\le\; 2\,D(F).
\]
As $\eta\to 0$, the bracket converges to $2 - \tfrac{2}{4} = \tfrac32$, so for $\delta_0$ small enough (forcing $\eta\le\eta_0(\eps)$ via Lemma~\ref{lem:NS}) we obtain
\[
  \B_0(F) \;\le\; \tfrac{2}{3/2 - \eps}\,D(F),
\]
which approaches $\tfrac{4}{3}\,D(F)$ as $\eps\to 0$. Setting $C_2 := 2/(3/2 - \eps)$ proves the claim.
\end{proof}

\subsection{Planar Closure Theorem}
\begin{theorem}[Planar Closure] \label{thm:Planar}
There exists a universal $C_2<\infty$ such that for every $2$-symmetric finite-perimeter set $F\subset\R^2$ with $|F|=\pi$,
\[ \beta(F)^2 \;\le\; C_2\,D(F). \]
\end{theorem}
\begin{proof}
Choose $\delta_0\in(0,1)$ from Lemma~\ref{lem:Fug}. If $D(F)\le\delta_0$, Lemma~\ref{lem:Fug} applies. If $D(F)>\delta_0$, Lemma~\ref{lem:Coarse} gives $\beta(F)^2 \le (1+2\pi/\delta_0)\,D(F)$. Taking the maximum of the two constants completes the proof.
\end{proof}

\begin{remark}[On the gap-1 Hardy obstruction]
The profile-graph approach, parametrizing $F$ by $|y|\le r(x)$, leads to the weighted inequality
\[
  \int_0^1 \rho\,\eta^2\,dx \;\le\; C\!\int_0^1 \rho^3 (\eta')^2\,dx, \qquad \eta=r-\rho,\ \rho=\sqrt{1-x^2},
\]
near the polar caps $x\to 1$. Substituting $s=1-x$ gives weights $s^{1/2}$ versus $s^{3/2}$ --- a \emph{gap of one} in exponent, whereas the standard Hardy inequality requires gap two. Test functions $\tilde\eta(s)=s^\beta$, $\beta\to 0^+$, refute the unconstrained gap-one estimate. The polar parametrization eliminates the problem at source: the equator (where the profile is regular) and the cap (where it is singular) are treated on equal footing as a single point on $S^1$, and the relevant Poincar\'e inequality on the compact circle has no boundary behavior to control.
\end{remark}

\section{Higher-Dimensional Symmetric Closure}

The polar Fuglede argument of Section~5 extends to dimension $n \ge 3$ verbatim, replacing the spectral gap of $-\partial_\theta^2$ on $S^1$ by that of the spherical Laplacian $-\Delta_{S^{n-1}}$. The $n$-symmetry of $F$ kills the constant mode (by volume) and all degree-one spherical harmonics (by reflection symmetry), leaving a spectral gap $\lambda_2 = 2n$ on the admissible subspace. No transverse Schwarz symmetrization, no dimensional induction, and no slice-wise potential-vs-perimeter-loss estimate is required.

\subsection{Coarse Bound at Large Deficit}
\begin{lemma}[Trivial Tail, dim $n$] \label{lem:CoarseN}
For any finite-perimeter $F\subset\R^n$ with $|F|=\omega_n$,
\[ \beta(F)^2 \le \B_0(F) \le n\omega_n + D(F). \]
Hence $\beta(F)^2 \le (1 + n\omega_n/\delta_0)\,D(F)$ whenever $D(F)\ge\delta_0$.
\end{lemma}
\begin{proof}
$\beta(F)^2 \le \B_0(F) = P(F) - (n-1)I_0(F) \le P(F) = n\omega_n + D(F)$, since $I_0(F)\ge 0$.
\end{proof}

\subsection{Nearly-Spherical Polar Parametrization}
\begin{lemma}[Fuglede Parametrization, dim $n$] \label{lem:NS_n}
There exist $\delta_0\in(0,1)$ and $C_*<\infty$ (universal in $n$) such that the following holds. If $F$ is $n$-symmetric (invariant under each reflection $x_i\to -x_i$), $|F|=\omega_n$, and $D(F)\le\delta_0$, then $F$ is star-shaped with respect to the origin and admits the polar parametrization
\[
   \partial F = \bigl\{(1+u(\xi))\,\xi : \xi \in S^{n-1}\bigr\}, \qquad u(R_i\xi) = u(\xi)\ \text{for each reflection } R_i,
\]
with
\[
   \|u\|_{L^\infty(S^{n-1})} + \|\nabla_{S^{n-1}} u\|_{L^\infty(S^{n-1})} \;\le\; C_*\,D(F)^{1/(4n)}.
\]
\end{lemma}
\begin{proof}[Proof sketch]
The Fusco--Maggi--Pratelli quantitative isoperimetric inequality gives an optimal centre $z_*$ with $|F\triangle B_1(z_*)| \le C_n\sqrt{D(F)}$; the $n$-symmetry forces $z_*=0$. The $L^1$-to-$W^{1,\infty}$ upgrade follows from $\Lambda$-minimality of the deficit-minimizing competitor and standard elliptic regularity (cf.\ Fuglede 1989; Fusco--Maggi--Pratelli, \emph{Ann.\ of Math.}\ 2008, §3). The 2-symmetry of $u$ is inherited from that of $F$. Star-shapedness follows from $\|u\|_\infty < 1$.
\end{proof}

\subsection{Polar Identity for $\B_0$}
\begin{lemma}[Polar Identity, dim $n$] \label{lem:PolarId_n}
For $F$ as in Lemma~\ref{lem:NS_n},
\[
   \B_0(F) - D(F) \;=\; \int_{S^{n-1}}\!\bigl[1 - (1+u(\xi))^{n-1}\bigr]\,d\sigma(\xi),
\]
and, to second order in $u$,
\[
   \B_0(F) \;=\; D(F) + \tfrac{n-1}{2}\!\int_{S^{n-1}}\! u^2\,d\sigma + R[u], \qquad |R[u]| \le C_n\,\|u\|_\infty\!\int_{S^{n-1}}\! u^2\,d\sigma.
\]
\end{lemma}
\begin{proof}
In polar coordinates with $\partial F = \{(1+u)\xi\}$:
\begin{align*}
   |F| &= \tfrac{1}{n}\!\int_{S^{n-1}}\!(1+u)^n\,d\sigma = \omega_n, \\
   I_0(F) &= \int_{S^{n-1}}\!\int_0^{1+u}\!\!s^{n-2}\,ds\,d\sigma = \tfrac{1}{n-1}\!\int_{S^{n-1}}\!(1+u)^{n-1}\,d\sigma, \\
   P(F) &= \int_{S^{n-1}}\!(1+u)^{n-2}\sqrt{(1+u)^2+|\nabla_{S^{n-1}} u|^2}\,d\sigma.
\end{align*}
Using $|S^{n-1}|=n\omega_n$ and $P(B_1)=n\omega_n$:
\[
  \B_0(F) - D(F) = n\omega_n - (n-1)I_0(F) = \int_{S^{n-1}}\!\bigl[1 - (1+u)^{n-1}\bigr]\,d\sigma.
\]
Expanding $(1+u)^k = 1 + ku + \binom{k}{2}u^2 + O(\|u\|_\infty u^2)$ for $k = n, n-1$, and substituting into the volume identity $\int[(1+u)^n - 1]d\sigma = 0$ to eliminate $\int u$ in favor of $-\tfrac{n-1}{2}\int u^2 + O(\|u\|_\infty\int u^2)$, the algebraic identity
\[
   n(1+u)^{n-1} - (n-1)(1+u)^n = 1 - \tfrac{n(n-1)}{2}u^2 + O(\|u\|_\infty u^2)
\]
gives $\int(1+u)^{n-1}d\sigma = n\omega_n - \tfrac{n-1}{2}\int u^2 + R$, hence $\B_0(F) - D(F) = \tfrac{n-1}{2}\int u^2 + R$ with the stated bound. \qedhere
\end{proof}

\subsection{Polar Fuglede Inequality in Dimension $n$}
\begin{lemma}[Polar Fuglede, dim $n$] \label{lem:Fug_n}
There exist $\delta_0\in(0,1)$ and $C_n<\infty$ such that for every $n$-symmetric $F$ with $|F|=\omega_n$ and $D(F)\le\delta_0$,
\[
   \beta(F)^2 \;\le\; \B_0(F) \;\le\; C_n\,D(F),
\]
with $C_n$ arbitrarily close to $\dfrac{2n}{n+1}$ for $\delta_0$ small.
\end{lemma}
\begin{proof}
By Lemma~\ref{lem:PolarId_n} (and absorbing $R[u]$ into $\eps$-corrections via $\|u\|_\infty \le C_* D^{1/(4n)} \to 0$), the claim reduces to
\[
   \int_{S^{n-1}}\!u^2\,d\sigma \;\le\; \tfrac{2(C_n - 1)}{n-1}\,D(F).
\]

\smallskip
\emph{Step 1 (Poincar\'e on $S^{n-1}$).} Decompose $u = c + \bar u$ with $c = (n\omega_n)^{-1}\int u\,d\sigma$ and $\int \bar u = 0$. Expand $\bar u$ in real spherical harmonics $\{Y_\ell^m\}_{\ell\ge 1}$. Each $Y_1^m$ is proportional to a coordinate $\xi_i$, which is odd under the reflection $R_i$; the $n$-symmetry of $\bar u$ kills every $\ell=1$ component. Hence $\bar u$ is supported on $\ell\ge 2$, where $-\Delta_{S^{n-1}}$ has eigenvalues $\ell(\ell+n-2) \ge 2n$, so
\[
   \int_{S^{n-1}}\!\bar u^2\,d\sigma \;\le\; \tfrac{1}{2n}\!\int_{S^{n-1}}\!|\nabla_{S^{n-1}}\bar u|^2\,d\sigma. \tag{$P_n$}
\]

\smallskip
\emph{Step 2 ($\B_0$ controls $\int|\nabla u|^2$).} The identity $\sqrt{a^2+b^2} - a = b^2/(\sqrt{a^2+b^2}+a)$ with $a=1+u$, $b=|\nabla u|$ gives
\[
   \B_0(F) \;=\; \int_{S^{n-1}} (1+u)^{n-2}\,\frac{|\nabla u|^2}{\sqrt{(1+u)^2+|\nabla u|^2} + (1+u)}\,d\sigma.
\]
For $\|u\|_\infty + \|\nabla u\|_\infty \le \eta \le 1/2$, the prefactor is bounded below by $(1-\eta)^{n-2}$ and the denominator above by $2(1+\eta) + \eta$. Hence
\[
   \int_{S^{n-1}}|\nabla u|^2\,d\sigma \;\le\; \frac{2+3\eta}{(1-\eta)^{n-2}}\,\B_0(F). \tag{$U_n$}
\]

\smallskip
\emph{Step 3 (control of the constant mode).} The volume identity $\int[(1+u)^n - 1]d\sigma = 0$ gives, to leading order, $|c| \le C_n\|u\|_\infty^2$ (the linear term in $c$ cancels against the quadratic $\int u^2$ term), hence
\[
   n\omega_n\,c^2 \;\le\; C_n\,\eta^2 \!\int_{S^{n-1}}\!\bar u^2\,d\sigma. \tag{$C_n$}
\]

\smallskip
\emph{Step 4 (Combine).} From $(P_n)$, $(U_n)$, $(C_n)$:
\[
   \int u^2 = \int\bar u^2 + n\omega_n c^2 \le (1 + C_n\eta^2)\!\int\bar u^2 \le \frac{(1 + C_n\eta^2)(2+3\eta)}{2n(1-\eta)^{n-2}}\,\B_0(F).
\]
By Lemma~\ref{lem:PolarId_n}, $\tfrac{2}{n-1}(\B_0 - D) = \int u^2 + O(\eta\int u^2)$, so
\[
   \B_0(F)\Bigl(\tfrac{2}{n-1} - \tfrac{(1 + C_n\eta^2)(2+3\eta)}{2n(1-\eta)^{n-2}} + O(\eta)\Bigr) \;\le\; \tfrac{2}{n-1}\,D(F).
\]
As $\eta\to 0$, the bracket converges to $\tfrac{2}{n-1} - \tfrac{1}{n} = \tfrac{n+1}{n(n-1)}$, giving
\[
   \B_0(F) \;\le\; \tfrac{2n}{n+1}\,D(F) \cdot (1 + o(1))_{\eta\to 0}.
\]
Setting $\delta_0$ small enough so that $\eta\le\eta_0$ (via Lemma~\ref{lem:NS_n}) absorbs the $o(1)$ into any chosen $C_n > 2n/(n+1)$.
\end{proof}

\begin{theorem}[Higher-Dimensional Symmetric Closure] \label{thm:HighDim}
For each $n\ge 2$ there exists $C_n<\infty$ such that for every $n$-symmetric finite-perimeter set $F\subset\R^n$ with $|F|=\omega_n$,
\[ \beta(F)^2 \;\le\; C_n\,D(F). \]
\end{theorem}
\begin{proof}
Choose $\delta_0$ from Lemma~\ref{lem:Fug_n}. If $D(F)\le\delta_0$, apply Lemma~\ref{lem:Fug_n}. Otherwise, Lemma~\ref{lem:CoarseN} gives $\beta(F)^2 \le (1+n\omega_n/\delta_0)D(F)$. Take the maximum of the two constants.
\end{proof}

\subsection{Combining with the Reflection Theorem}

Putting Theorem~\ref{thm:OFR} (oriented reflection to an $n$-symmetric $F$) together with Theorem~\ref{thm:HighDim} (Fuglede closure for the $n$-symmetric $F$) yields the main result.

\begin{theorem}[Main Theorem: Sharp $\B$-Oscillation--Deficit Inequality] \label{thm:Main}
For every finite-perimeter $E\subset\R^n$ with $|E|=\omega_n$,
\[ \beta(E)^2 \;\le\; \tilde C_n\,D(E). \]
\end{theorem}
\begin{proof}
By Theorem~\ref{thm:OFR}, $\beta(E)^2 \le C_n(D(E) + \beta(F)^2)$ for an $n$-symmetric $F$ produced by iterated median reflection. By Theorem~\ref{thm:HighDim}, $\beta(F)^2 \le C_n D(F)$. Each reflection step preserves the deficit up to a constant (the median-cut perimeter inequality $P(A_1^\pm) \le P(A) + 2P(A;H_i) \le 3P(A)$, hence $D(A_1^\pm) \le C D(A)$), so $D(F) \le C_n^{\,n} D(E)$. Combining,
\[ \beta(E)^2 \le C_n(D(E) + C_n \cdot C_n^{\,n} D(E)) = \tilde C_n\,D(E). \qedhere\]
\end{proof}

\begin{remark}[Why polar, not profile-graph]
The profile-graph approach $|y|\le r(x)$ uses coordinates that are degenerate at the polar caps $x=\pm 1$: the Hessian of the perimeter functional has rank $1$ rather than $2$, and the resulting weighted Hardy estimate has gap one rather than gap two. The polar parametrization $\xi\mapsto (1+u(\xi))\xi$ on $S^{n-1}$ trades the cap singularity for a compact-domain Poincar\'e estimate with the genuine spherical-Laplacian spectral gap. The $n$-symmetry of $F$ then knocks out exactly the $\ell=1$ translation modes that would otherwise prevent $\B_0 \le C_n D$.
\end{remark}

\end{document}
